
\chapter{Background and related work}
\label{ch:background}

A conversational memory system determines which parts of an interaction history remain
available to a language model and in what form. Keeping the original dialogue, retrieving
selected passages and constructing a graph make different demands on the model and on the
device running it. This chapter separates the information an application stores from the
resources needed to use it, introduces the memory mechanisms evaluated in the thesis, and
reviews the work most relevant to their accuracy and resource costs on small models and
smartphones.

Much of the closest work is recent and was available only as an arXiv preprint when this chapter
was written. This includes the two on-device systems SmartRAG and MobileRAG, the source-expansion
system LeanMem, and the memory systems LazyMem, MemFlow, MemoryCPT, RippleMem and
MOSAIC~\citep{smartrag,mobilerag_energy,leanmem,lazymem,memflow,memorycpt,ripplemem,mosaic}.
Preprints have not been peer reviewed. They are cited here for the designs and measurements their
authors report. They give context for the thesis's results but are not evidence for them. Every reference
was checked against its arXiv or publisher record. Where a preprint has since
appeared at a conference or in a journal, the bibliography cites that version.

\section{Small language models and local inference}
\label{sec:ch2_slms}

\subsection{The model within an agent}

A language model predicts tokens conditioned on the sequence supplied to it. An agent
places that model inside a system that can receive observations, call tools, maintain
state and act on its environment. Memory is one component of this system: it supplies
information from earlier interactions that is relevant to a current decision. The survey
by \citet{memory_survey2024} distinguishes the model's capabilities from the mechanisms
that store, retrieve and process its interaction history.

The argument for using small language models (SLMs) in agents rests partly on this division
of work. An application may repeatedly need a restricted operation, such as classifying a
request or answering from supplied evidence. \citet{slm_agentic} argue that compact models
can serve such operations economically and that agent systems can combine models of
different sizes. This is a design argument for selecting models by their assigned task;
it does not establish that a small model can perform every memory operation reliably.
Extracting relations, resolving conflicting statements and answering from retrieved text
remain distinct capabilities that require evaluation.

There is no single parameter threshold that determines whether a model is practical on a
phone. Architecture, weight precision, context length and runtime support also affect its
resource requirements. The on-device review by \citet{ondevice_lm_review} therefore treats
model design, compression and hardware execution together. In an empirical study of
compact models, \citet{slm_edge_deployment} likewise find that model size alone does not
explain capability or execution speed. This thesis uses SLM to describe its evaluated
compact readers; Chapter~\ref{ch:methodology} specifies their identities and parameter
scales. As in Chapter~\ref{ch:introduction}, the reader is the model that receives the
selected evidence and generates the answer. When another model prepares the memory, its
role and cost must be specified separately.

\subsection{Context windows and persistent memory}

Three meanings of memory need to be distinguished. \emph{Parametric memory} refers to
information encoded in model weights through training. The \emph{context window} contains
the tokens available to a particular inference call. \emph{External memory} is information
stored by the surrounding application and made available to later calls. Retrieval-augmented
generation connects an external store to a language model without requiring each
new fact to be incorporated into its weights~\citep{rag_lewis}. MemGPT makes the distinction
between limited in-context working space and longer-term external storage explicit through
a hierarchy of memory tiers~\citep{memgpt}.

Appending an interaction to a prompt does not, by itself, update the model's weights. If
that interaction is omitted from a later call, the application needs another mechanism to
recover it. A large context window postpones the point at which history must be selected
or compressed, but it does not specify how a growing history should be managed. Moreover,
input capacity and effective use of evidence are different properties. \citet{lost_in_middle}
show that answer quality can depend on where relevant information appears within a long
input, including for models designed to accept long contexts. Their result motivates
testing context use rather than treating the advertised context length as a guarantee of
recall.

In a bounded full-context baseline, the reader receives as much of the history as the
configured input limit permits. Once that limit is exceeded, a truncation policy determines
which material remains visible. Such a baseline is a useful reference for the cost of
reading dialogue directly, but its name must not imply that every historical turn is
present. Retrieval addresses a different selection problem: it chooses evidence according
to the current query instead of relying only on the portion of the history that fits.

\subsection{Prefill, decoding and quantization}

Autoregressive inference has two stages with different resource demands. During
\emph{prefill}, the model processes the input tokens and prepares the internal state used
for generation. During \emph{decoding}, it generates subsequent tokens sequentially. For
attention layers, cached key and value tensors allow the model to reuse representations of
earlier tokens. This \emph{KV cache} occupies runtime memory in addition to the model
weights. Its size depends on the retained sequence length and the attention architecture,
so a model whose weights fit in RAM can still become impractical with a large
context~\citep{ondevice_lm_review,slm_edge_deployment}.

Weight quantization reduces the precision used to represent model parameters. It can
reduce storage and memory-bandwidth demand, but its benefits depend on the implementation and
hardware. Weight precision also does not determine the precision or size of every runtime
buffer. In their edge measurements, \citet{slm_edge_deployment} distinguish the predominantly
compute-bound prefill stage from the predominantly memory-bandwidth-bound decode stage and
show that quantization affects them differently. The relevant bottleneck can therefore
change when a memory strategy replaces a long input with additional generated summaries
or extraction outputs.

This distinction matters for cost accounting. Reader input tokens, reader output tokens
and tokens generated while constructing memory perform different jobs. Counting them is
useful, but adding them into one total conceals the operations responsible for that cost.
A resource study must record both the work performed and the device resources consumed.

\section{Representations of long-term conversational memory}
\label{sec:ch2_memory}

\subsection{The memory life cycle}

External memory involves decisions at several points in an interaction. At ingestion, the
system chooses what to retain and may transform raw observations into summaries, facts or
structured records. During maintenance, it may merge related entries, revise outdated
information or remove material. At query time, retrieval selects candidate evidence and
context assembly decides how to present it to the reader. The agent-memory and graph-memory
surveys organize these operations as connected parts of a memory life
cycle~\citep{memory_survey2024,graph_memory_survey2026}.

Conversational information also differs in how it changes. A stable preference may be
useful across many sessions, while an appointment has a date and a decision may later be
reversed. Preserving only the latest statement can lose the evidence needed to answer a
historical question. THEANINE addresses this issue by linking memories through temporal
and causal relations and retrieving timelines, retaining earlier information that helps
explain a change~\citep{theanine}. Its design illustrates why forgetting and updating
policies must be chosen in relation to the task.

Table~\ref{tab:ch2_memory_representations} summarizes four recurring representation choices.
These are overlapping design families. A graph can retain pointers to raw turns, and a
vector index can store summaries instead of verbatim text. The three experimental arms
in this thesis select particular configurations from this broader space;
Figure~\ref{fig:ch3_lookup_methods} in Chapter~\ref{ch:methodology} shows how each of the three
looks up evidence.

\begin{table}[htbp]
  \centering
  \small
  \begin{tabular}{p{0.18\textwidth}p{0.34\textwidth}p{0.36\textwidth}}
    \toprule
    Representation & Information supplied to the reader & Main trade-off \\
    \midrule
    Full context & Original history within an input limit & Preserves included wording;
    input grows with history and truncation can remove evidence. \\
    Vector retrieval & Text records selected by embedding similarity & Limits reader input;
    retrieval may miss supporting records or return redundant text. \\
    Compressed memory & Summaries or selected factual records & Reduces repeated text;
    construction costs work and may discard future query details. \\
    Graph memory & Retrieved entities, relations, events or linked source passages & Supports
    associative access; usefulness depends on construction, traversal and representation. \\
    \bottomrule
  \end{tabular}
  \caption{Memory representation choices, synthesized from the reviewed literature
  \citep{memory_survey2024,rag_lewis,comedy,graph_memory_survey2026}. The trade-offs describe
  design considerations, not a measured ranking.}
  \label{tab:ch2_memory_representations}
\end{table}

\subsection{Dense retrieval over dialogue}

Dense retrieval represents a query and stored text as numerical vectors. DPR establishes
a dual-encoder approach in which questions and passages are mapped into a shared space
and ranked using a similarity score~\citep{dpr}. Contriever demonstrates that contrastive
training can produce useful dense retrieval without supervised relevance labels for each
target dataset~\citep{contriever}. These methods provide foundations for indexing dialogue
with an existing embedding model, although performance on conversational histories still
depends on the encoder and the way records are formed.

For a query $q$ and stored record $d_i$, a common score is cosine similarity:
\begin{equation}
  s(q,d_i)=\frac{f(q)^{\mathsf T}g(d_i)}{\lVert f(q)\rVert\,\lVert g(d_i)\rVert},
  \label{eq:ch2_similarity}
\end{equation}
where $f$ and $g$ denote the query and record encodings. The top-$k$ records are candidates
for the reader's context. This notation describes the retrieval operation used in the
thesis; dense retrievers can use other scoring functions, including the unnormalized
inner product used by DPR.

The retrieval unit affects what a successful match contains. A short turn can isolate one
fact but leave a pronoun unresolved. A longer passage can retain its surrounding exchange
while also introducing irrelevant content. Speaker identifiers and timestamps can be
necessary for interpreting the retrieved words. These are representation choices before
the reader sees the evidence, not properties of the reader itself.

Increasing $k$ can bring more supporting material into the candidate set, but it also
increases the amount of text that must be selected, ordered and read. Chain-of-Memory
examines the resulting separation between retrieving relevant evidence and answering
correctly from it~\citep{chain_of_memory}. Consequently, retrieval coverage and answer
accuracy are complementary evaluation targets. High coverage identifies available
evidence; it does not establish that the context assembled from that evidence is useful
to a particular reader.

\subsection{Compression, hierarchy and memory management}

These systems share one decision: how much interpretation to perform before the future
question is known. Some compress what is stored. COMEDY maintains a compressive memory in
place of a retrieval database~\citep{comedy}, RAPTOR recursively clusters and summarizes text
into a tree of abstraction levels~\citep{raptor}, and MemoRAG uses a global memory of the input
to generate retrieval clues~\citep{memorag}. Others manage records as they arrive. MemGPT lets
the model move information between in-context and external storage~\citep{memgpt}, Mem0
extracts and updates salient facts, with a graph variant for relations~\citep{mem0}, and A-MEM
links notes and lets new memories revise existing ones~\citep{amem}. Hierarchical designs add
structure on top: \citet{hmem_yu2026} pair a temporal and semantic tree with a knowledge graph
of the entities it mentions, and the H-MEM system of \citet{hmem_sun2025} routes retrieval
between levels of semantic abstraction, with LoCoMo scores that belong to its own reader and
protocol.

Early compression can reduce later reader work, but it commits to a view of what is worth
preserving. Retaining raw dialogue keeps more information available and moves selection work
to query time, and systems that keep source pointers beside compact records partly defer the
decision. This thesis approaches the same question from the other side: what a triple loses
when the text behind it is not shown.

\section{Graph retrieval and preservation of evidence}
\label{sec:ch2_graphs}

\subsection{What the graph represents}

A knowledge graph represents entities and relations between them. A subject--relation--object
triple, such as $(\text{person},\text{works at},\text{organization})$, is one possible unit.
Conversational systems may attach timestamps, confidence values and source references to
that unit. Other graph memories connect complete events or text passages instead of
individual entities. The graph-memory taxonomy of \citet{graph_memory_survey2026} covers
these different representations and the operations used to construct and retrieve them.

GraphRAG is therefore a family label as well as the name of a specific system. The
approach of \citet{graphrag_edge2024} builds an entity graph, forms communities and generates
community summaries for query-focused summarization over a corpus. HippoRAG combines
knowledge-graph construction with Personalized PageRank to retrieve information through
associative connections~\citep{hipporag2024}. LightRAG uses graph and vector representations
with retrieval at both detailed entity and broader thematic levels~\citep{light_rag}.
Different choices of graph nodes, edges and retrieval procedures lead to different
evidence being shown to the reader.

The graph arm of this thesis uses SLM-extracted triples from dialogue and retrieves
evidence from that store. It is a particular graph-memory pipeline, described in
Chapter~\ref{ch:methodology}. Its behavior cannot stand for every system bearing the
GraphRAG name. In particular, a community-summary system and a graph linking intact
episodes preserve and expose information differently from a triple-only prompt.

\subsection{Construction, retrieval and reader limitations}

Graph construction can introduce errors before retrieval begins. The extractor must
identify the relevant entities and relations, preserve their meaning and produce an
output that the storage pipeline can interpret. Valid syntax alone does not establish
that a fact was extracted correctly. A missing relation and a relation that has lost a
qualification are different errors, and both can affect later answering. The loss can be
large when a small model performs the extraction. In a multi-hop kinship-reasoning task,
\citet{kelesis2026} report that an RGCN model reaches 60.69\% accuracy on ground-truth
facts but 24.88\% on facts extracted by a 4B Gemma model.

The controlled comparison by \citet{rag_vs_graphrag} shows that graph results depend on the
construction model, the retrieval design and the form of evidence supplied to the reader,
including triples alone and triples with source text. Giving plain retrieval the same token
budget as community-based graph retrieval on MultiHop-RAG removes most of the graph's
advantage: with a 70B reader, token-matched retrieval reaches 71.01\% against 71.17\% for the
graph, compared with 65.77\% under the default budget, although the graph stays ahead on temporal
and comparison questions. Representation and context volume are therefore necessary controls in a
graph comparison. Preparation cost also differs: the same study reports 135 seconds to build its
RAG index against 5,560 to 7,702 seconds for its graph variants.

Conversational systems with richer graphs show that graph memory can work well. RippleMem
connects event-centred memory units through participants, locations and time, retrieving
anchors and expanding through their associations~\citep{ripplemem}. MOSAIC stores entity-typed
nodes, checks new items for contradictions at save time, and replaces an LLM classification step
with locality-sensitive hashing at negligible accuracy cost~\citep{mosaic}. Both build richer
memory units than triples emitted by a local small model, and their reported LoCoMo results
reflect those designs. A graph result is therefore a property of the whole construction and
retrieval design, and the expensive part of that design need not be an additional language-model
call.

\subsection{Provenance and source-text expansion}

\emph{Provenance} records where a memory item came from. A pointer from a triple or summary
to the original dialogue can support inspection and allow the reader to recover details
that the compact representation omitted. For example, a statement about a former employer
and a statement about a current employer can collapse into the same undated relation if
temporal wording is discarded. A source pointer makes the original distinction available
without requiring it to be reproduced in every compact record.

SEEM combines an episodic event layer with a graph layer and connects both to their source
passages. Its Reverse Provenance Expansion retrieves original passages through the
provenance links of recalled event frames~\citep{seem}. LeanMem distinguishes stable profile
memory, evolving event memory and detail-preserving record memory. Its records retain
pointers to original dialogue spans, which can be expanded when a query requires detailed
evidence~\citep{leanmem}. Both systems explicitly address the loss of information when
interaction history is condensed.

SEEM and LeanMem are therefore the closest related work for the Selective Provenance
Expansion (SPE) study in Chapter~\ref{ch:spe}. The specific question examined there is how
to choose among omitting a retrieved fact, retaining its triple and supplying its source
turn under a fixed retrieval-context budget. The study evaluates a loss-aware selection
rule with compact readers; its scope is the behavior of that rule and its implementation.

\section{Efficient construction and budget-aware memory use}
\label{sec:ch2_budgets}

\subsection{When to construct memory}

An efficient memory design must account for work performed before a question is asked. The AI
Hippocampus finds on LongMemEval-S that a simple vector store can be competitive with more
elaborate memory systems while needing much less processing time~\citep{ai_hippocampus}, so the
benefit of construction must be measured against a plain retrieval baseline. Other systems change
when and how processing happens. Chain-of-Memory keeps preparation light and organizes retrieved
fragments into memory chains at query time~\citep{chain_of_memory}. LazyMem and DeferMem retrieve a
broad candidate pool and use a trained lightweight model to compress or rewrite it into
query-conditioned evidence~\citep{lazymem,defermem}. MemoryCPT instead distils the construction
pipeline into one compact model and scores systems by quality per unit of inference cost, measured
as a dollar price over language-model tokens~\citep{memorycpt}; a distilled construction model is
one way to reduce the phone construction energy measured in Chapter~\ref{ch:on_device_memory}.

Deferral changes the cost distribution. It reduces eager, query-independent construction but adds
processing to every query, while storing or indexing observations still costs work. An amortization
calculation for a graph built once and reused across many questions therefore does not describe
every memory architecture; a query-time construction system needs its own measurement boundary,
including any training or offline preparation relevant to the deployment.

\subsection{Routing and selective processing}

Budget-aware memory systems decide how much work a query warrants. BudgetMem offers
memory-processing modules at low, medium and high budget tiers, varying implementation
complexity, reasoning effort or model capacity, and trains a compact router among
them~\citep{budgetmem}; Chapter~\ref{ch:memory_evaluation} places the retrieval-depth sweep
relative to these axes. MemFlow addresses compact readers directly: intent routing chooses among
profile lookup, targeted retrieval and deeper reasoning, evidence is assembled under a tiered
token budget, and a validator can escalate to a heavier path~\citep{memflow}. MemSkill learns
which memory-processing skills to apply and revises them from difficult cases~\citep{memskill}.
Outside memory, Selective Context removes less informative input before
inference~\citep{selective_context}, and TALE conditions reasoning on a token
budget~\citep{tale}.

These systems measure cost in tokens, calls, latency or price, and separate execution cost from
the offline work used to obtain a policy. The reviewed evaluations do not convert their decisions
into measured phone energy, so their savings cannot be read as energy saved on a particular
handset.

\section{On-device memory and physical resource costs}
\label{sec:ch2_device}

\subsection{The closest mobile studies}

On-device execution brings memory design into contact with shared RAM, battery capacity,
processor scheduling and heat dissipation. \citet{slm_edge_deployment} evaluate compact
models on edge hardware and smartphones and observe differences in prompt processing,
decoding and thermal behavior. Their study supports treating runtime and hardware as part
of the experimental condition. The on-device survey by \citet{ondevice_lm_review} also
describes hardware acceleration and model compression as complementary approaches to
deployment.

\citet{mobile_memory_architectures} compare raw buffers, tiered memory, vector retrieval
and structured graphs on short- and long-term mobile-agent applications. Their findings
show that the preferred architecture depends on the application and model capacity, with
structured memory remaining competitive in their tested settings. This is a close
conceptual comparison because it varies memory architecture with compact models. Its
latency measurements use an M2 Pro laptop, and its graph construction uses deterministic
rules. These choices differ from phone execution with an SLM extractor. Its MemGUI-Bench
and FingerTip tasks also differ from the conversational-QA benchmarks used here. The
authors list real-phone latency and device-level cost measurement as future work;
Chapter~\ref{ch:on_device_memory} reports both for conversational memory.

MemLoRA targets the memory operations themselves. It distills specialized adapters for
extraction, updating and response generation into compact models and evaluates their use
in a conversational memory system~\citep{memlora}. This provides an alternative to relying
on an unadapted small extractor: memory-specific training can change the quality of the
intermediate representation. The paper's throughput measurements, however, do not provide
a physical smartphone energy comparison. It addresses memory capability and computational
efficiency under its reported execution setup.

Two further systems treat memory as part of a phone-native agent runtime. MOBIMEM reports
CPU-only experiments on a Snapdragon 8 Elite platform and finds that action reuse can
reduce latency by factors of 1.6$\times$ to 9$\times$, without reporting handset
energy~\citep{mobimem}. ClawMobile treats token budgets, inference placement and persistent
memory as smartphone-runtime concerns, but its preliminary agent uses remote model
inference~\citep{clawmobile}.

MobileRAG provides direct device-energy evidence for document retrieval. It combines
EcoVector, a vector search design that partially loads index data, with Selective Content
Reduction to reduce the text passed to the reader. Its evaluation on a Galaxy S24 includes
energy per query for document-QA workloads~\citep{mobilerag_energy}. Its focus is document
indexing and query processing rather than representations of long-term conversation history.

SmartRAG is the closest reviewed on-device graph system. It combines a distilled,
continually adaptable entity recognizer and relation extraction with a provenance-preserving
graph, hybrid retrieval and a compact answer model. It runs on two OnePlus handsets with
1.7B and 4B backbones and reports memory use and average power during extraction and
QA~\citep{smartrag}. These measurements show that a graph-based retrieval system can be
deployed on a phone when its components are designed for that setting.

Two boundaries distinguish its evidence from the comparison needed here. SmartRAG's QA
benchmarks use document corpora, not long-term dialogue histories. In addition,
Appendix~C states that its competing retrieval systems run in a cloud-side simulation;
they are not profiled as complete real-phone deployments. Its average-power measurements
therefore characterize the proposed system rather than establish a matched per-answer
energy comparison across memory strategies~\citep{smartrag}.

\subsection{Tokens, time, power and energy}

Cost measures answer different questions. Token counts describe the volume of model input
and output. Latency describes the elapsed time for an operation. Power, measured in watts,
is the rate of energy use; energy, measured in joules, accumulates over the operation's
duration. For a measurement interval of length $T$,
\begin{equation}
  E=\int_0^T P(t)\,\mathrm{d}t.
  \label{eq:ch2_energy}
\end{equation}
An average-power value can be converted to energy only with the duration for the same
measurement interval. A lower power draw can coexist with higher task energy when an
operation takes sufficiently longer.

The difference between MobileRAG's per-query energy and SmartRAG's average power illustrates
why units and boundaries must remain explicit~\citep{mobilerag_energy,smartrag}. Similarly,
tokens generated by an extractor and tokens read by an answer model cannot be assigned a
universal energy coefficient. Prefill and decoding have different bottlenecks, and model
architecture and hardware affect their execution~\citep{slm_edge_deployment}. Token counts
can be informative predictors within a measured configuration, but the relationship must
be established for that configuration.

Memory construction also changes the time horizon of an efficiency claim. For a fixed
store used by $N$ queries, one accounting model is
\begin{equation}
  \overline{E}(N)=\frac{E_{\mathrm{build}}+E_{\mathrm{maint}}}{N}
  +\frac{1}{N}\sum_{j=1}^{N}E_{\mathrm{query},j},
  \label{eq:ch2_amortization}
\end{equation}
where $E_{\mathrm{build}}$ is construction energy, $E_{\mathrm{maint}}$ is maintenance
energy over the same period, and each query term includes the operations inside the chosen
service boundary. This is an accounting identity for the stated workload, not a claim
that costs stay constant as a store grows. Query-time construction, as in LazyMem or
DeferMem, belongs in the query term~\citep{lazymem,defermem}.

Accuracy and energy per query describe answer quality and device cost separately; a combined
measure asks how much correct output the energy buys. \citet{intelligence_per_watt} compare local
inference configurations by accuracy per watt and accuracy per joule. Because power omits
duration (Equation~\ref{eq:ch2_energy}), only the per-joule form separates a strategy that answers
in one minute from one that needs several minutes at a similar draw. Correct answers per
kilojoule, used in this thesis, is the per-workload form of accuracy per joule. It describes a
device workload only when its correctness labels come from the answers that workload produced, and
it complements rather than replaces accuracy, energy, latency, RAM and thermal measurements.
Chapter~\ref{ch:methodology} gives its operational definition.

\section{Evaluating conversational memory}
\label{sec:ch2_evaluation}

\subsection{LoCoMo and LongMemEval}

LoCoMo evaluates memory over long, multi-session conversations. The published benchmark
includes question answering, event summarization and multimodal dialogue generation,
with conversations averaging approximately 600 turns and 16,000 tokens and spanning up to
32 sessions~\citep{locomo}. Its QA component includes questions requiring individual
facts, evidence across turns, temporal interpretation and open-domain knowledge, together
with an adversarial category. This thesis evaluates the QA component using the specified
first-four-category convention. Its results do not assess the benchmark's summarization
or multimodal-generation tasks.

LongMemEval tests information extraction, multi-session reasoning, temporal reasoning,
knowledge updates and abstention through 500 curated questions embedded in interaction
histories~\citep{longmemeval}. These abilities separate recalling an isolated fact from
combining sessions, identifying the relevant point in time or recognizing that the history
does not support an answer. Its inclusion provides a second conversational-memory setting
for the thesis. The actual subset and evaluation configuration are defined in
Chapter~\ref{ch:methodology}.

The benchmarks are related but are not interchangeable. In particular, abstention is an
explicit LongMemEval ability, while the thesis's primary LoCoMo analysis excludes its
adversarial category under the stated convention. Coverage of one question type cannot
be inferred from success on another. Both benchmarks also assess memory through specified
tasks; performance on their QA questions is evidence about memory-dependent answering,
not a direct measure of tool use or complete mobile-agent workflows.

\subsection{Automatic judging and comparison boundaries}

Short-answer metrics based on lexical overlap can miss semantically correct paraphrases
or respond to differences in answer wording. An LLM judge can instead assess a candidate
answer against the question and reference answer using a fixed rubric. This makes the
judge part of the measurement procedure. MT-Bench and Chatbot Arena study model-based
judging and document biases, including sensitivity to answer position and verbosity,
alongside agreement with human preferences in their setting~\citep{mtbench}. Agreement
reported for those tasks cannot calibrate a different judge on conversational-memory QA.

Repeatability and validity must consequently be evaluated separately. Repeated agreement
by one judge indicates stability under the tested procedure. Human review addresses
whether its labels are appropriate for the task. A reference-guided, single-answer
protocol also differs from pairwise preference judging, so biases established for one
format should not be assumed to transfer unchanged. LongMemEval already uses the
reference-guided format: it scores answers with an expert-written correctness prompt and
reports more than 97\% agreement between its GPT-4o judge and human
experts~\citep{longmemeval}. That figure validates the prompt with that judge model; it does
not transfer automatically to a different model applying the same rubric, which is the
arrangement used in this thesis. The methodology chapter specifies
the judging procedure and the limits of the available validation.

The same care is needed when comparing published systems. Two studies using LoCoMo may
use different conversations, exclude different question categories, train on part of the
corpus, allow different answer lengths, or use different readers and judges. DeferMem,
for example, separately reports LoCoMo evaluation excluding the conversation used to
train its distiller~\citep{defermem}. That result belongs to a different evaluation set
from a study that retains all conversations. The benchmark name alone is insufficient
to establish a matched comparison.

For this reason, the related work is compared by method, task and measurement boundary.
Published accuracy values establish what a system achieved under its own protocol.
The thesis's controlled comparisons use paired questions and a common scoring procedure
within the experimental conditions described in Chapter~\ref{ch:methodology}.

\section{Synthesis and position of the thesis}
\label{sec:ch2_position}

The reviewed work establishes several viable approaches to long-term memory. Dense
retrieval preserves access to original text, while compression and graph construction
change how information is organized and exposed. Recent systems improve the use of
retrieved evidence, defer construction until the query is known, or restore source text
through provenance. Their results make memory design a question of representation,
processing and evaluation conditions rather than a fixed ordering of architectures.

Table~\ref{tab:ch2_related_position} summarizes the closest connections. It records the
type of evidence supplied by each group of studies, keeping system feasibility, proxy
efficiency and physical energy measurements distinct.

\begin{table}[htbp]
  \centering
  \small
  \begin{tabular}{p{0.23\textwidth}p{0.30\textwidth}p{0.35\textwidth}}
    \toprule
    Related work & Relevant contribution & Relation to this thesis \\
    \midrule
    MobileRAG~\citep{mobilerag_energy} & Smartphone document RAG with per-query energy &
    Direct energy precedent; different stored content and comparison scope. \\
    SmartRAG~\citep{smartrag} & On-phone graph retrieval with power and memory profiles &
    Graph feasibility precedent; competing systems lack matched phone profiles. \\
    Mobile memory benchmark~\citep{mobile_memory_architectures} & Architecture comparison
    for compact mobile agents & Different tasks; laptop latency measurements and
    rule-based graph construction. \\
    MemLoRA~\citep{memlora} & Distilled adapters for small-model memory operations &
    Alternative construction design; no matched handset-energy study. \\
    RippleMem; Chain-of-Memory~\citep{ripplemem,chain_of_memory} & Associative retrieval
    and deliberate evidence assembly & Richer memory units performing well; limits how far a
    triple-pipeline result generalizes. \\
    LazyMem; DeferMem~\citep{lazymem,defermem} & Query-conditioned construction after broad
    retrieval & Alternative distribution of preparation and query costs. \\
    BudgetMem; MemFlow~\citep{budgetmem,memflow} & Query-aware choice of memory processing
    and budgets & Established routing approaches; phone energy remains a separate measure. \\
    LeanMem; SEEM~\citep{leanmem,seem} & Compact memory connected to original source text &
    Direct related work for the budgeted provenance-selection study. \\
    \bottomrule
  \end{tabular}
  \caption{Closest related work and its measurement boundaries. Entries summarize
  the versions reviewed for this chapter and do not form a cross-paper leaderboard.}
  \label{tab:ch2_related_position}
\end{table}

Read across its rows, the table separates three kinds of evidence: feasibility on a phone,
efficiency measured by proxies such as tokens, latency or price, and physical energy. Among the
studies reviewed here, per-query phone energy is reported for document
retrieval~\citep{mobilerag_energy}, and average phone power for one graph system whose baselines
ran off the phone~\citep{smartrag}. The closest comparison of memory architectures for compact
mobile agents reports latency on a laptop and leaves real-phone cost measurement to future
work~\citep{mobile_memory_architectures}. These studies do not compare memory representations for
conversational histories by answer quality and energy measured on the same handset. This thesis
makes that comparison for the tested memory choices and analyses where its graph pipeline loses
evidence; the following chapter defines the experimental controls needed to interpret it.
